% ════════════════════════════════════════════════════════════════
%  前言（frontmatter，置于目录与正文之间；罗马页码）
%  说明：草稿，作者可改写。措辞不绑定书名，便于后续定名。
% ════════════════════════════════════════════════════════════════
\chapter*{前言}
\addcontentsline{toc}{chapter}{前言}
\markboth{前言}{前言}

本书脱胎于作者学习机器人学时的个人笔记。这一领域已有若干杰出的著作：有的以几何直觉与工程代码见长，有的以严谨的状态估计理论取胜，有的则以现代记号与系统广度著称。然而，能把它们的视角熔于一炉、又把每一步推导都摊开讲透的读物，并不多见。本书的初衷，正是做这样一次“全量吸收、多源合并、逐步推导”的整理——把分散在不同著作与论文中的精华，按\emph{主题}（而非按文献）重组为自洽的一体。

全书以机器人学的数学与理论为骨架——刚体运动、李群与李代数、状态估计与滤波、非线性优化、SLAM 与多传感器感知、运动规划与控制——并循序扩展到\emph{具身智能}：强化学习、视觉–语言–动作（VLA）模型，以及大模型驱动的策略与操作。换言之，本书试图沿“经典模型 $\to$ 学习方法 $\to$ 具身智能”的脉络，架起一座由传统机器人学通向现代具身智能的桥梁。

在取材与写法上，本书坚持几条原则。其一，\textbf{自包含}：所引著作与论文中有价值的内容，均在书中\emph{完整写出}，每一步推导不跳步，引文仅用于标注出处，绝不以“详见原书”代替正文。其二，\textbf{多视角合并}：同一概念若有不同视角，各自独立复述再融会贯通。其三，\textbf{逐公式溯源}：关键结论标注来源，便于回查与延伸。

为求全书一致，若干记号与约定贯穿始终：流形上的求导与更新以\textbf{右扰动}为主线（局部坐标，便于对接因子图与 GTSAM/Ceres），四元数采用 \textbf{Hamilton} 约定（实部在前），$\se(3)$ 坐标取“平移在前” $\boldsymbol{\xi}=[\boldsymbol{\rho};\boldsymbol{\phi}]$。完整的符号与记号约定见随后的《主要符号与记号约定》一章，建议先行浏览。

各理论章大体遵循“前置自测 $\to$ 学习目标 $\to$ 知识导航 $\to$ 各节（动机 $\to$ 反面 $\to$ 理论 $\to$ 陷阱 $\to$ 练习）$\to$ 常见误解 $\to$ 小结 $\to$ 故障排查”的脉络，难度以 $\bigstar$ 标注（导论、综述、结语等非理论章不拘此式）；书中以不同的彩色盒区分内容类型——定义/定理/命题、推导、代码、洞见、陷阱与练习。读者既可顺序通读，亦可按需查阅。

本书是持续生长的工程，内容会不断增补修订，记号亦可能随之微调；后续版本将在作者个人仓库 \url{https://github.com/Michael-Jetson/Robotics_Tutorial} 发布（仓库地址暂定）。疏漏谬误在所难免，恳请读者不吝指正。

\vspace{1em}
\noindent\textbf{致谢}\quad 本书是一次“站在巨人肩上”的整理：各主题的骨架与精华，皆得益于下列著作的滋养。谨向这些作者致以诚挚谢意——没有它们，便没有本书。

\begin{itemize}\setlength{\itemsep}{2pt}
  \item \textbf{视觉 SLAM 与状态估计}：高翔等《视觉 SLAM 十四讲：从理论到实践》\cite{gaoxiang2019slam14}；T.~D.~Barfoot, \emph{State Estimation for Robotics}\cite{barfoot2024state}；L.~Carlone 等编《The SLAM Handbook》\cite{carlone2026handbook}；F.~Dellaert 与 M.~Kaess, \emph{Factor Graphs for Robot Perception}\cite{dellaert2017factor}。
  \item \textbf{李群与随机模型}：G.~S.~Chirikjian, \emph{Stochastic Models, Information Theory, and Lie Groups}\cite{chirikjian2011stochastic}。
  \item \textbf{运动规划与控制}：S.~M.~LaValle, \emph{Planning Algorithms}\cite{lavalle2006planning}；K.~J.~{\AA}str{\"o}m 与 R.~M.~Murray, \emph{Feedback Systems}\cite{astrom2008feedback}；H.~K.~Khalil, \emph{Nonlinear Systems}\cite{khalil2002nonlinear}；J.~B.~Rawlings 等, \emph{Model Predictive Control}\cite{rawlings2020mpc}；F.~Borrelli 等, \emph{Predictive Control for Linear and Hybrid Systems}\cite{borrelli2017predictive}；M.~W.~Spong 等, \emph{Robot Modeling and Control}\cite{spong2006robot}；R.~Tedrake, \emph{Underactuated Robotics}\cite{tedrake_underactuated}。
  \item \textbf{四足与多刚体动力学}：卞泽坤、王兴兴（杭州宇树科技）《四足机器人控制算法》\cite{unitree2023quadruped}；于宪元《基于稳定性的仿生四足机器人控制系统设计》\cite{yu2021quadruped}；R.~Featherstone, \emph{Rigid Body Dynamics Algorithms}\cite{featherstone2008rigid}。
  \item \textbf{强化学习}：赵世钰《强化学习的数学原理》\cite{zhao2024rlmath}；R.~S.~Sutton 与 A.~G.~Barto, \emph{Reinforcement Learning: An Introduction}\cite{sutton2018rl}。
\end{itemize}

\noindent 此外，正文中以引用标注的大量论文、教程与开源实现，其作者的贡献同样不可或缺，恕不一一列举，在此一并致谢。文中若有错漏，责任全在整理者，与上述著作无关。

\vspace{1.5em}
\begin{flushright}
  gpf\\[0.3em]
  2026 年
\end{flushright}
